% ============================================================
% f2014_chongqing_science_q05_flowchart.tex
% 2014 重庆卷（理科）第 5 题：程序框图
% 来源 PDF：普通高考/2014/2014重庆理.pdf
% 原页：第 1 页（A3 横排）左栏，印面 865，第 5 题题干下方
% 结构：开始 → 赋初值 k=9、s=1 → 判断框（框内待填条件）；
%       是：s=s·k/(k+1) → k=k-1 → 返回判断；否：输出 k → 结束。
% ============================================================
\begin{tikzpicture}[>=Stealth, line width=0.7pt, font=\small,
  line cap=round, line join=round]

  % 开始
  \draw[rounded corners=6pt] (-0.95,0.45) rectangle (0.95,-0.45);
  \node at (0,0) {开始};
  \draw[->] (0,-0.45) -- (0,-0.95);

  % 初始化
  \draw (-1.75,-0.95) rectangle (1.75,-1.85);
  \node at (0,-1.40) {$k=9$，$s=1$};

  % 进入判断框
  \draw[->] (0,-1.85) -- (0,-2.85);

  % 判断菱形（原页框内为空，条件待填）
  \draw (0,-2.85) -- (2.15,-3.90) -- (0,-4.95) -- (-2.15,-3.90) -- cycle;

  % 否：向下输出
  \draw[->] (0,-4.95) -- (0,-5.45);
  \node at (-0.38,-5.20) {否};
  \draw (-1.95,-5.45) -- (1.55,-5.45) -- (1.95,-6.45) -- (-1.55,-6.45) -- cycle;
  \node at (0,-5.95) {输出 $k$};
  \draw[->] (0,-6.45) -- (0,-6.95);

  % 结束
  \draw[rounded corners=6pt] (-0.95,-6.95) rectangle (0.95,-7.85);
  \node at (0,-7.40) {结束};

  % 是：向右进入循环体
  \draw[->] (2.15,-3.90) -- (2.95,-3.90);
  \node at (2.55,-3.58) {是};

  % 循环体第一框
  \draw (2.95,-4.40) rectangle (6.45,-3.40);
  \node at (4.70,-3.90) {$s=s\cdot\dfrac{k}{k+1}$};

  % 循环体第二框
  \draw[->] (4.70,-3.40) -- (4.70,-2.70);
  \draw (2.95,-2.70) rectangle (6.45,-1.70);
  \node at (4.70,-2.20) {$k=k-1$};

  % 回边：向左汇入进入判断框的竖线
  \draw (2.95,-2.20) -- (0,-2.20);
\end{tikzpicture}
